\documentclass[11pt,a4paper]{article}
\usepackage[T1]{fontenc}
\usepackage{lmodern,amsmath,amssymb,amsthm,mathtools,booktabs,microtype}
\usepackage[margin=27mm]{geometry}
\usepackage[colorlinks=true,linkcolor=blue,urlcolor=blue,citecolor=blue]{hyperref}
\newtheorem{theorem}{Theorem}
\newtheorem{proposition}[theorem]{Proposition}
\newtheorem{lemma}[theorem]{Lemma}
\newtheorem{corollary}[theorem]{Corollary}
\newtheorem{definition}[theorem]{Definition}
\newcommand{\Q}{\mathbb Q}\newcommand{\Z}{\mathbb Z}
\newcommand{\A}{\mathbb A}\newcommand{\PP}{\mathbb P}
\DeclareMathOperator{\rad}{rad}
\DeclareMathOperator{\sfpart}{sf}
\setlength{\emergencystretch}{2em}
\title{Additional arithmetic hypotheses\\and polynomial definitions of the integers}
\author{}\date{9 November 2026}
\begin{document}\maketitle
\begin{abstract}
We distinguish height inequalities that shorten a polynomial definition
from prime-existence statements and from conjectures about rational
points. A weak $abc$ inequality with exponent $121/100$ suffices for
the thirteen-coordinate construction; consequently the ordinary
discriminant--conductor form of Szpiro's conjecture suffices.
The truncated Vojta conjecture supplies the stronger ordinary $abc$
inequality. We prove precise limitations of two other proposed
reductions: distinct prime assignments do not control the squarefree
part of a product, and a polynomial injection does not by itself
give a finite algebraic encoding of several free coordinates.
Finally, a topological argument explains the incompatibility between
Mazur's connected-component conjecture and an existential definition
of the integers. No reduction in the number of coordinates from the
Riemann hypothesis or Grimm's conjecture is asserted.
\end{abstract}

\section{Definitions and the required height estimates}

All coordinate counts include the input. A polynomial
$F\in\Z[X,Y_1,\ldots,Y_m]$ defines $\Z$ in $\Q$ if
\[
 x\in\Z\quad\Longleftrightarrow\quad
 \exists\mathbf y\in\Q^m\quad F(x,\mathbf y)=0.
\]
For a nonzero integer $n$, put
\[
 \rad(n)=\prod_{p\mid n}p,\qquad
 \sfpart(n)=\prod_{\substack{p\mid n\\v_p(n)\ {\rm odd}}}p.
\]
For $t=u/v$ in lowest terms with $v>0$, let
$h(t)=\log\max(|u|,v)$. If $(w,u,v)$ is a primitive common
integer representative of $(1,a,b)$ with $w>0$, define
$h(1:a:b)=\log\max(w,|u|,|v|)$.

\begin{definition}
For $\kappa>1$, the weak $abc$ inequality with exponent $\kappa$
is the assertion that a constant $C_\kappa>0$ exists such that
\[
 \max(|a|,|b|,|c|)
 \le C_\kappa\rad(abc)^\kappa
\]
for every primitive nonzero integer triple $a+b=c$.
\end{definition}

The ordinary rational $abc$ conjecture asserts this for every
$\kappa>1$. The following result uses only one exponent.

\begin{theorem}\label{thm:weak}
Assume the weak $abc$ inequality with exponent $121/100$ and the
elliptic group, pole obstruction and realization, and simultaneous
prime-pattern inputs specified in~\cite{Thirteen}.
There is a computable coefficient family in thirteen total
coordinates whose sufficiently large fixed odd input shifts define
$\Z$ in $\Q$. A sufficient numerical shift is not determined.
The same conclusion follows under the ordinary Szpiro conjecture
and the stated inherited inputs.
\end{theorem}
\begin{proof}
We describe the numerical inequalities that make the extra
hypothesis sufficient; the complete rational source and polynomial
conversion are given in~\cite{Thirteen}. The rational map
\[
 f(t)=\frac{t(t+2)^3}{(2t+1)^3},
 \qquad
 f(t)-1=\frac{(t-1)^3(t+1)}{(2t+1)^3}
\]
has degree four. Its inverse images of $0,1,\infty$ are the
six rational points $0,-2,-1/2,1,-1,\infty$.
Applying the weak inequality to its numerator, denominator, and
their difference gives
\[
 \log M(t)\ge(8/\kappa-6)h(t)-O(1),
\]
where $M$ is the product of primes having odd contact valuation,
outside a fixed exceptional set. To see the coefficient, the
radical of the associated degree-six binary form is at least a
constant times $\exp(4h(t)/\kappa)$, whereas
$\rad(n)^2\le |n|\sfpart(n)$ and the form has size
$O(\exp(6h(t)))$. The fixed change $t=2+1/\lambda$ makes all
six contacts finite.

The same weak inequality, applied to coprime eighth powers, gives
\[
 h(1:\beta:T)\le\frac{121}{437}h(\beta^8+T^8)+O(1).
 \tag{1}
\]
The common-denominator calculation proving (1) is included
in~\cite{Thirteen}; its general coefficient is
$(d/\kappa-3)^{-1}$ for even $d>3\kappa$.
For the two arguments $\lambda$ and $(\beta^8+T^8)\lambda$,
the contact bound and the local congruences therefore give
\[
 h(1:\beta:T)\le
 \frac{14641}{16169}\log N+O(1)
 \tag{2}
\]
whenever $N=|\operatorname{num}(\beta-n)|\ne0$.
Here $\beta=2X^2+B$, the elliptic index $n$ is an integer,
and the one-point height theorem gives $n^2\le A h(T)+J$
with fixed constants. Since
$\log N\le h(\beta)+\log(1+|n|)$, (2) bounds both $n$ and
$h(\beta)$ independently of the fixed odd shift $B$.
Choosing $B$ larger than the exponential of the bound for
$h(\beta)$ forces $\beta=n$.
The elementary denominator calculation for $2X^2+B$ then gives
$X\in\Z$. Completeness and the thirteen-coordinate count follow
from the explicit source in~\cite{Thirteen}.

The semistable case of the ordinary Szpiro conjecture is the assertion
$|\Delta_E|\le C_\epsilon N_E^{6+\epsilon}$
for semistable elliptic curves over $\Q$, with minimal
discriminant $\Delta_E$ and conductor $N_E$.
Its implication of the weak $abc$ inequality with every exponent
$6/5+\eta$ is established in~\cite[\S I.3]{Oesterle}.
Take $\eta=1/100$.
\end{proof}

This implication uses the discriminant--conductor conjecture itself.
It does not identify that conjecture with the stronger modified
Szpiro statement about $\max(|c_4|^3,|c_6|^2)$.
The modified statement is equivalent to ordinary $abc$;
see~\cite[\S I.3]{Oesterle}. The thirteen-coordinate conclusion
does not remove the pole or elliptic realization hypotheses.

\section{The truncated Vojta conjecture}

Let $X$ be smooth and projective over $\Q$, let $D$ be a
reduced divisor with normal crossings, and let $A$ be ample.
The rational-point case of the truncated Vojta conjecture
asserts, outside a proper closed exceptional subset, that
\[
 N^{(1)}(D,P)\ge h_{K_X+D}(P)-\epsilon h_A(P)-O(1).
 \tag{3}
\]
The counting function records a prime once when the reduction of
$P$ meets $D$. This is the specialization of the conjecture
in~\cite[Conjecture 2.3]{Vojta}.

\begin{proposition}
The case of (3) for
$X=\PP^1$, $D=[0]+[1]+[\infty]$, and $A=[\infty]$
implies the ordinary rational $abc$ conjecture.
\end{proposition}
\begin{proof}
Write $P=a/c$ with $a+b=c$ primitive and nonzero.
At finite primes the truncated counting function is
$\log\rad(abc)$. The divisor $K_X+D$ has degree one, so
its height is $h(P)+O(1)$. Thus (3) gives
\[
 (1-\epsilon)h(P)\le\log\rad(abc)+O(1).
\]
The difference between $h(P)$ and
$\log\max(|a|,|b|,|c|)$ is bounded by $\log 2$.
A proper closed subset of $\PP^1$ is finite, so its remaining
rational points can be absorbed into the constant.
Taking $\epsilon=\eta/(1+\eta)$ gives exponent $1+\eta$.
\end{proof}

Consequently this Vojta hypothesis supplies every additional
$abc$ input of the thirteen-coordinate and conditional quartic
constructions. A lower coordinate count does not follow merely
from replacing $abc$ by this stronger conjecture.

\section{Prime assignments and odd valuations}

Grimm's conjecture asserts that a finite block of consecutive
composite positive integers admits distinct prime divisors,
one for each member. Its formulation and its relation to prime
gaps are discussed in~\cite{ErdosPomerance}.
The next elementary example isolates a difficulty in using
such assignments as a height estimate.

\begin{proposition}\label{prop:pell}
There are arbitrarily large positive integers $n$ such that
$n$ and $n+1$ are both composite and
\[
                    \sfpart(n(n+1))=2.
\]
Each such pair admits distinct prime divisors.
\end{proposition}
\begin{proof}
Define positive integers $a_k,b_k$ by
\[
 (17+6\sqrt8)^k=a_k+b_k\sqrt8,\qquad k\ge1.
\]
Their norm identity is $a_k^2-8b_k^2=1$, and $a_k$ is odd.
For $n_k=(a_k-1)/2$ this gives
$n_k(n_k+1)=2b_k^2$.
Multiplication by $17+6\sqrt8$ acts by the integral matrix
\[
 \begin{pmatrix}17&48\\6&17\end{pmatrix},
\]
whose determinant is one. Its image in
$\mathrm{GL}_2(\Z/144\Z)$ has some finite order $L$.
For $k\equiv1\pmod L$ we therefore have
$a_k\equiv17\pmod{144}$, whence $n_k\equiv8\pmod{72}$.
In particular $8\mid n_k$ and $9\mid n_k+1$; both are
composite. These $n_k$ tend to infinity. Their consecutive
members are coprime, so any prime divisor of one differs
from every prime divisor of the other.
\end{proof}

Proposition~\ref{prop:pell} concerns the information supplied
by a prime assignment. It is not an independence theorem
between conjectures. In particular, it shows why assigning
different primes to factors cannot be substituted for a
lower bound on primes occurring to odd order.

The Riemann hypothesis concerns the zeros of the Riemann zeta
function; its generalized form concerns zeros of the appropriate
$L$-functions. Their usual prime-counting consequences concern
the distribution of primes, as illustrated by
Schoenfeld's conditional estimates~\cite{Schoenfeld}.
Such a counting estimate does not supply the uniform height
inequalities (1)--(2). No implication from either hypothesis
to those inequalities, or to a smaller defining polynomial,
is used here. Similarly, Schinzel-type prime-value hypotheses
could alter a completeness argument only after the relevant
simultaneous polynomial values and admissibility conditions
have been specified. They do not by themselves replace the
soundness argument.

\section{Polynomial injections and finite algebraic encodings}

The Bombieri--Lang conjecture predicts that the rational points
of a projective variety of general type are not Zariski dense.
For surfaces, Poonen proves that this conjecture implies the
existence of a polynomial inducing an injection
$\Q^2\to\Q$~\cite[Theorem 1.1]{Poonen}.
His Remark 1.7 records a different sufficient hypothesis:
the four-term generalized $abc$ conjecture.
It is essential to distinguish that conjecture from ordinary,
three-term $abc$.

\begin{proposition}\label{prop:dimension}
Let $k$ be a characteristic-zero field and let $U$ be a nonempty
open subset of $\A^r_k$, with $r>1$.
No rational function $j:U\dashrightarrow\A^1_k$ is generically
finite onto its image.
\end{proposition}
\begin{proof}
If $j$ is constant, its only fiber has dimension $r$.
Otherwise $k(j)$ has transcendence degree one over $k$, whereas
$k(U)=k(t_1,\ldots,t_r)$ has transcendence degree $r$.
Thus
$\operatorname{trdeg}_{k(j)}k(U)=r-1>0$,
which is incompatible with a finite function-field extension.
\end{proof}

A polynomial may nevertheless be injective on $k$-rational
points, since geometric and rational fibers are different.
If $j$ is injective on $\Q^r$ and a relation is
$F(x,\mathbf y)=0$, encoding its rational witnesses by
$z=j(\mathbf y)$ gives the exact reformulation
\[
 \exists\mathbf y\in\Q^r:
   F(x,\mathbf y)=0,\quad z=j(\mathbf y).
 \tag{4}
\]
Injectivity proves uniqueness of a rational inverse when one
exists; it does not supply its existence or remove the
quantifiers in (4).
Proposition~\ref{prop:dimension} also excludes replacing that
inverse on an open set by a finite algebraic extension, the
setting required by the one-coordinate separation theorem.
This is a limitation of that proposed encoding. It does not
rule out a different use of polynomial injections in a
Diophantine construction.

\section{A conjecture incompatible with the projection}

Mazur's connected-component conjecture asserts that for every
variety $V/\Q$, the closure of $V(\Q)$ in $V(\mathbb R)$ has
finitely many connected components. The implication below is
classical; see~\cite[\S\S 2--3]{CornelissenZahidi}.

\begin{proposition}
If Mazur's connected-component conjecture holds, $\Z$ is
not Diophantine in $\Q$.
\end{proposition}
\begin{proof}
Suppose a polynomial $F(x,\mathbf y)$ defines $\Z$, and let
$V$ be its affine zero set. Let $C$ be the real topological
closure of $V(\Q)$ in $V(\mathbb R)$. The coordinate projection
$\pi:V(\mathbb R)\to\mathbb R$ is continuous. Since
$\pi(V(\Q))=\Z$ and $\Z$ is closed,
$\pi(C)\subseteq\Z$. Conversely every integer occurs at a
rational point, so $\pi(C)=\Z$.
Every connected component of $C$ has a connected image in
the discrete set $\Z$, hence a one-point image.
Finitely many components would give a finite image, a
contradiction. If the conjecture is formulated only for
irreducible varieties, apply it to the finitely many
irreducible components of $V$.
\end{proof}

Thus this topological conjecture cannot be adjoined to
arithmetic premises already implying an existential definition
of $\Z$ as a consistent method of reducing the number of witnesses.
It is distinct from the Bombieri--Lang statement used above.

\begin{thebibliography}{99}
\bibitem{Thirteen}
\emph{Thirteen variables from a weak $abc$ inequality},
companion paper and exact construction,
\url{../conditional_few_variables/conditional13_20261109_v16.pdf}.
\bibitem{Oesterle}
J.~Oesterl\'e,
Nouvelles approches du \emph{th\'eor\`eme} de Fermat,
S\'eminaire Bourbaki, Exp.~694,
\emph{Ast\'erisque} \textbf{161--162} (1988), 165--186.
\url{https://www.numdam.org/article/SB_1987-1988__30__165_0.pdf}.
\bibitem{Vojta}
P.~Vojta, A more general $abc$ conjecture,
\emph{Int. Math. Res. Notices} (1998), no.~21, 1103--1116.
\url{https://arxiv.org/abs/math/9806171}.
\bibitem{ErdosPomerance}
P.~Erd\H{o}s and C.~Pomerance,
An analogue of Grimm's problem of finding distinct prime
factors of consecutive integers,
\emph{Utilitas Math.} \textbf{24} (1983), 45--65.
\url{https://math.dartmouth.edu/~carlp/PDF/38.pdf}.
\bibitem{Schoenfeld}
L.~Schoenfeld, Sharper bounds for the Chebyshev functions
$\theta(x)$ and $\psi(x)$. II,
\emph{Math. Comp.} \textbf{30} (1976), 337--360.
\url{https://doi.org/10.1090/S0025-5718-1976-0457374-X}.
\bibitem{Poonen}
B.~Poonen, Multivariable polynomial injections on rational numbers,
\emph{Acta Arith.} \textbf{145} (2010), 123--127.
\url{https://math.mit.edu/~poonen/papers/QQ.pdf}.
\bibitem{CornelissenZahidi}
G.~Cornelissen and K.~Zahidi,
Topology of Diophantine sets: remarks on Mazur's conjectures,
in \emph{Hilbert's Tenth Problem: Relations with Arithmetic
and Algebraic Geometry}, Contemp. Math.~270, 2000, 253--260.
\url{https://arxiv.org/abs/math/0006140}.
\end{thebibliography}
\end{document}
